%% chapters/chapter02/chapter02.tex
%% CHAPTER 2: LITERATURE REVIEW

\chapterfigpath{chapter02}

\chapter{Literature Review}
\label{chap:literature_review}


\section{Overview}
\label{sec:overview}

This chapter reviews the existing literature related to the research topic. The review is organized thematically to provide a comprehensive understanding of the current state of knowledge.

Figure~\ref{fig:literature_taxonomy} presents a taxonomy of the literature reviewed in this chapter.

\begin{figure}[ht]
    \centering
    \includegraphics[width=0.75\textwidth]{literature_taxonomy}
    \caption{Taxonomy of literature reviewed showing main themes and their relationships}
    \label{fig:literature_taxonomy}
\end{figure}


\section{Theoretical Foundations}
\label{sec:theoretical_foundations}

Discuss the theoretical foundations relevant to your research. This section should cover fundamental concepts, theories, and models.

A fundamental equation in this field is:

\begin{equation}
    f(x) = ax^2 + bx + c
    \label{eq:example_equation}
\end{equation}

where $a$, $b$, and $c$ are constants, and $x$ is the input variable.


\section{Existing Approaches}
\label{sec:existing_approaches}

Table~\ref{tab:approach_comparison} compares different approaches found in the literature.

\begin{table}[ht]
    \centering
    \caption{Comparison of existing approaches}
    \label{tab:approach_comparison}
    \begin{tabular}{@{}lccc@{}}
        \toprule
        \textbf{Approach} & \textbf{Advantage} & \textbf{Limitation} & \textbf{Reference} \\
        \midrule
        Method A & Fast execution & Lower accuracy & \cite{sample_ref1} \\
        Method B & High accuracy & Slower speed & \cite{sample_ref2} \\
        Method C & Balanced & Complex setup & \cite{sample_conf1} \\
        \bottomrule
    \end{tabular}
\end{table}


\section{Technologies and Tools}
\label{sec:technologies}

Discuss the technologies, frameworks, or tools relevant to your research. This section provides context for the technical decisions made in the methodology chapter.


\section{Evaluation Methods}
\label{sec:evaluation_methods}

Algorithm~\ref{alg:example} shows a sample algorithm structure commonly used in this domain.

\begin{algorithm}[H]
\caption{Example Algorithm}
\label{alg:example}
\begin{algorithmic}[1]
\Require Input data $D$, Parameter $k$
\Ensure Output result $R$
\State Initialize $R \gets \emptyset$
\For{each item $i$ in $D$}
    \If{condition is met}
        \State Process item $i$
        \State Add result to $R$
    \EndIf
\EndFor
\State \Return $R$
\end{algorithmic}
\end{algorithm}

\section{Related Work}
\label{sec:related_work}

Several studies have addressed similar problems:

\textbf{Study by Author et al. \cite{sample_ref1}:} Briefly describe what they did and their findings. Then explain how your work differs or extends this.

\textbf{Study by Another Author \cite{sample_ref2}:} Another similar study and its relevance to your work.

\textbf{Study by Third Author \cite{sample_conf1}:} Continue with more related work.


\section{Research Gap}
\label{sec:research_gap}

While existing literature provides extensive coverage of this topic, there are gaps in:

\begin{enumerate}
    \item Gap 1 - What hasn't been done that you will do
    \item Gap 2 - What limitation in existing work you will address
    \item Gap 3 - What new perspective or approach you will bring
\end{enumerate}

This research addresses these gaps by providing a comprehensive solution.


\section{Chapter Summary}
\label{sec:lit_review_summary}

This chapter reviewed the main areas of literature relevant to this research. Key findings include the identification of existing approaches, their limitations, and the research gap that this study aims to address. The identified research gaps motivate the methodology presented in the following chapter.
